% !TeX root = Report.tex
\chapter{Численное сравнение методов предобработки перед преобразованием Хафа}

\section{Дизайн численного эксперимента}

В данной главе сравниваются методы предобработки изображения перед применением
преобразования Хафа. Рассматриваются методы \texttt{CSSA ROW.ROW}, \texttt{CSSA COL.ROW},
\texttt{MAX.ROW}, Median и Wiener. Отдельно анализируется вариант, в котором
после CSSA применяется ESPRIT c коррекцией частот.
Для методов сравниваются три варианта формирования активного множества: без
порога, с порогом $1\hat\sigma$ и с порогом $2\hat\sigma$, где $\hat\sigma$
оценивается по остаточной матрице. Метод \texttt{MAX.ROW} устроен иначе: он
выбирает $k$ максимальных элементов в каждой строке, поэтому пороговая
фильтрация для него не является основным этапом метода.

В экспериментах используются два уровня аддитивного гауссовского шума:
$\sigma=0.05$ и $\sigma=0.20$. Для одной прямой рассматриваются четыре
конфигурации, то есть прямые $y=x$, $y=-x+101$, $y=2x-1$ и $y=2x-40$,
во всех случаях с интенсивностью $0.8$. Для двух прямых используется пара
$y=2x-1$ с интенсивностью $0.8$ и $y=-x+101$ с интенсивностью $0.5$.

Во всех таблицах ниже используются три метрики: MSE по $\rho$, MSE по $\theta$
и среднее число активных пикселей после предобработки. Чем меньше первые две
метрики, тем точнее восстановлены параметры прямой. Число активных пикселей
позволяет интерпретировать причины высоких ошибок: если после предобработки
остается слишком много или слишком мало активных пикселей, то в
аккумуляторном массиве могут возникать ложные максимумы, приводящие к большим
ошибкам в оценках $\rho$ и $\theta$.

Во всех методах на основе CSSA используется стандартная длина окна $L=50$, поскольку
обрабатываемые ряды имеют длину $100$. В варианте \texttt{ROW.ROW} CSSA применяется к строкам
матрицы после DFT, а в варианте \texttt{COL.ROW} --- к ее столбцам. Число
сохраняемых компонент выбирается автоматически по конфигурации прямых:
для одной прямой в \texttt{ROW.ROW} сохраняется одна компонента, а в
\texttt{COL.ROW} сохраняется одна компонента для прямых $y=x$ и $y=-x+101$
и две компоненты для прямых $y=2x-1$ и $y=2x-40$. Для двух прямых сохраняются
две компоненты в \texttt{ROW.ROW} и три компоненты в \texttt{COL.ROW}.
В варианте \texttt{CSSA ROW.ROW + ESPRIT} используется то же число компонент,
что и в \texttt{ROW.ROW}; параметры экспонент оцениваются методом ESPRIT,
после чего частоты округляются к сетке с шагом $1/N_{\text{col}}$. Метод
\texttt{MAX.ROW} выбирает $k=1$ максимум в каждой строке для одной прямой и
$k=2$ максимума для двух
прямых. Median использует окно $3\times3$, а Wiener --- локальное окно
$5\times5$. В вариантах с пороговой фильтрацией активными считаются элементы,
не меньшие $m\hat\sigma$, где $m\in\{1,2\}$. В варианте без пороговой
фильтрации активными считаются положительные элементы после предобработки.

Эксперименты для одной прямой проводятся на матрицах размера
$N_{\text{row}}\times N_{\text{col}}=100\times100$ при размере выборки
$n=100$ реализаций; параметры сведены в табл.~\ref{tab:ch3-exp-03}.
Используемая дискретизация в этом случае задается шагами
$h_\rho=0.05$ и $h_\theta=0.005$.

\begin{table}[H]
	\centering
	\caption{Параметры эксперимента для одной прямой}
	\label{tab:ch3-exp-03}
	\resizebox{0.85\textwidth}{!}{%
		\begin{tabular}{|c|c|c|c|c|c|c|}
\hline
$N_{\text{row}}$ & $N_{\text{col}}$ & $n$ & $\sigma$ & пороги & $h_\rho$ & $h_\theta$ \\
\hline
100 & 100 & 100 & $\{0.05,\,0.20\}$ & без порога, $1\hat\sigma$, $2\hat\sigma$ & 0.05 & 0.005 \\
\hline
\end{tabular}

	}
\end{table}

Для экспериментов с двумя прямыми также используется
выборка из $n=100$ матриц размера $100\times100$; соответствующие шаги
дискретизации составляют $h_\rho=0.02$ и $h_\theta=0.002$
(табл.~\ref{tab:ch3-exp-04}).

\begin{table}[H]
	\centering
	\caption{Параметры эксперимента для двух прямых}
	\label{tab:ch3-exp-04}
	\resizebox{0.85\textwidth}{!}{%
		\input{chapters/chapter3_experiments/tables/ch3_exp_params_04}
	}
\end{table}

\section{Влияние дискретизации преобразования Хафа}

Для интерпретации последующих таблиц сначала зафиксируем вклад дискретизации.
В табл.~\ref{tab:ch3-discretization-decay} приведен расчет для
прямой $y=2x-1$ при нулевом шуме: истинные параметры прямой ищутся по сетке,
после чего вычисляются MSE по $\rho$ и $\theta$ относительно истинных
параметров. В таблицу включены шаги сетки, используемые далее в
экспериментах.

\begin{table}[H]
	\centering
	\caption{Влияние шага сетки Хафа на нижний уровень ошибки для прямой
	$y=2x-1$}
	\label{tab:ch3-discretization-decay}
	\resizebox{0.55\textwidth}{!}{%
		\begin{tabular}{|c|c|c|c|}
\hline
$h_\rho$ & $h_\theta$ & MSE по $\rho$ & MSE по $\theta$
\\
\hline
5.0000 & 1.0000 & 0.200000 & 0.103719
\\
\hline
1.0000 & 0.0100 & 0.200000 & 4.22e-06
\\
\hline
0.5000 & 0.0050 & 0.002786 & 4.22e-06
\\
\hline
0.2500 & 0.0025 & 0.002786 & 1.98e-07
\\
\hline
0.0500 & 0.0050 & 7.76e-06 & 4.22e-06
\\
\hline
0.0200 & 0.0020 & 5.20e-05 & 3.02e-09
\\
\hline
\end{tabular}

	}
\end{table}

Конкретное значение зависит не только от шага, но и от положения истинных параметров
относительно узлов сетки, поэтому далее ошибки методов сравниваются с нижним
уровнем, соответствующим выбранной дискретизации.

\section{Без пороговой фильтрации}

\subsection{Одна прямая}

В данном разделе сравниваются методы предобработки в задаче обнаружения одной
прямой. При анализе результатов важно учитывать не только MSE оценок $\rho$ и
$\theta$, но и число активных элементов. Эта величина показывает, насколько
сильно предобработка меняет множество точек, по которому строится
аккумуляторный массив. Если после обработки остается слишком много активных
пикселей, то в пространстве параметров появляются ложные максимумы; если их
слишком мало, часть линии фактически теряется.

В табл.~\ref{tab:ch3_one_005_nothr} и
\ref{tab:ch3_one_020_nothr} приведены результаты для уровней шума
$\sigma=0.05$ и $\sigma=0.20$ соответственно.

\begin{table}[H]
\centering
\caption{Одна прямая: $\sigma=0.05$, без порога}
\label{tab:ch3_one_005_nothr}
\resizebox{\textwidth}{!}{%
\begin{tabular}{|l|l|c|c|c|c|c|c|}
\hline
Конфигурация & Метрика & \texttt{ROW.ROW} & \texttt{COL.ROW} & \texttt{MAX.ROW} & Median & Wiener & \texttt{ROW.ROW}+ESPRIT
\\
\hline
$(-1,101;c=0.8)$ & MSE по $\rho$ & 3.16e-04 & 79.7457 & 3.16e-04 & 1346.6 & 4957.3 & 3.16e-04
\\
\hline
 & MSE по $\theta$ & 1.59e-07 & 1.59e-07 & 1.59e-07 & 0.6045 & 0.6169 & 1.59e-07
\\
\hline
 & Активные пиксели & 1960.7 & 4536.5 & 100 & 4924.1 & 9999.8 & 100
\\
\hline
$(1,0;c=0.8)$ & MSE по $\rho$ & 4173.2 & 3017.8 & 2.50e-03 & 3621.8 & 1.1500 & 2.50e-03
\\
\hline
 & MSE по $\theta$ & 0.6181 & 0.6169 & 1.43e-06 & 0.6169 & 0.6169 & 1.43e-06
\\
\hline
 & Активные пиксели & 1968.6 & 4570.8 & 100 & 4932.6 & 9999.8 & 100
\\
\hline
$(2,-1;c=0.8)$ & MSE по $\rho$ & 3488.4 & 2622.1 & 0.0233 & 2568.1 & 0.3898 & 0.0233
\\
\hline
 & MSE по $\theta$ & 3.8613 & 0.2150 & 4.22e-06 & 0.2284 & 0.2150 & 4.22e-06
\\
\hline
 & Активные пиксели & 2093.9 & 4674.2 & 100 & 4864.1 & 9999.8 & 74.9
\\
\hline
$(2,-40;c=0.8)$ & MSE по $\rho$ & 2281.6 & 1774.5 & 0.0447 & 1704.2 & 282.0 & 0.0447
\\
\hline
 & MSE по $\theta$ & 3.9005 & 0.2150 & 4.22e-06 & 0.2150 & 0.2150 & 4.22e-06
\\
\hline
 & Активные пиксели & 2075.2 & 4655.7 & 100 & 4849.1 & 9999.7 & 75
\\
\hline
\end{tabular}%
}
\end{table}


\begin{table}[H]
\centering
\caption{Одна прямая: $\sigma=0.20$, без порога}
\label{tab:ch3_one_020_nothr}
\resizebox{\textwidth}{!}{%
\begin{tabular}{|l|l|c|c|c|c|c|c|}
\hline
Конфигурация & Метрика & \texttt{ROW.ROW} & \texttt{COL.ROW} & \texttt{MAX.ROW} & Median & Wiener & \texttt{ROW.ROW}+ESPRIT
\\
\hline
$(-1,101;c=0.8)$ & MSE по $\rho$ & 695.1 & 57.8912 & 3.16e-04 & 1570.8 & 4955.9 & 3.16e-04
\\
\hline
 & MSE по $\theta$ & 0.3324 & 1.59e-07 & 1.59e-07 & 0.6107 & 0.6169 & 1.59e-07
\\
\hline
 & Активные пиксели & 4062.8 & 4866.6 & 100 & 5040.2 & 10000 & 85.2
\\
\hline
$(1,0;c=0.8)$ & MSE по $\rho$ & 3351.7 & 3003.4 & 5.08e-03 & 3412.3 & 1.0300 & 7.23e-03
\\
\hline
 & MSE по $\theta$ & 0.6181 & 0.6169 & 1.43e-06 & 0.6169 & 0.6169 & 1.43e-06
\\
\hline
 & Активные пиксели & 4114.2 & 4927.1 & 100 & 5089.2 & 10000 & 84.9
\\
\hline
$(2,-1;c=0.8)$ & MSE по $\rho$ & 3106.1 & 2561.0 & 0.0200 & 2985.2 & 0.3266 & 0.0194
\\
\hline
 & MSE по $\theta$ & 3.7740 & 0.2284 & 4.22e-06 & 0.2150 & 0.2150 & 4.31e-06
\\
\hline
 & Активные пиксели & 4006.8 & 4980.1 & 100 & 5002.3 & 9999.9 & 66.6
\\
\hline
$(2,-40;c=0.8)$ & MSE по $\rho$ & 1443.8 & 1952.2 & 0.0317 & 1668.6 & 284.9 & 0.0250
\\
\hline
 & MSE по $\theta$ & 3.8221 & 0.2150 & 4.22e-06 & 0.2150 & 0.2150 & 4.27e-06
\\
\hline
 & Активные пиксели & 3961.1 & 4983.4 & 100 & 5004.5 & 10000 & 67.4
\\
\hline
\end{tabular}%
}
\end{table}


Для конфигурации
$(-1,101;c=0.8)$ метод \texttt{ROW.ROW} сохраняет малые ошибки:
$3.16\cdot 10^{-4}$ по $\rho$ и $1.59\cdot 10^{-7}$ по $\theta$ при
$1960.7$ активных пикселей. Однако при сопоставимом числе активных пикселей
для других конфигураций ошибки \texttt{ROW.ROW} значительно возрастают.
Например, для $(1,0;c=0.8)$ MSE по $\rho$ равно $4173.2$, а MSE по $\theta$
равно $0.6181$. Это различие показывает, что само по себе большое число
активных пикселей не является достаточным условием ошибки, но делает результат
преобразования Хафа чувствительным к конфигурации прямой.

Отдельно стоит отметить поведение фильтра Wiener. В некоторых наклонных
конфигурациях он дает сравнительно малую ошибку по $\rho$, например для
$y=2x-1$, но ошибка по $\theta$ остается большой. Причина в том, что без
порога после Wiener почти вся матрица остается активной, поэтому расстояние до
прямой может оцениваться приемлемо, тогда как направление выбирается по
побочному максимуму. Методы \texttt{MAX.ROW} и
\texttt{CSSA ROW.ROW + ESPRIT} в этих же конфигурациях сохраняют небольшое
число активных пикселей и поэтому не создают сопоставимого числа ложных
направлений в аккумуляторном массиве.

\subsection{Две прямые}

Для двух прямых используется та же схема анализа: после предобработки
формируется активное множество, по которому строится аккумуляторный массив
преобразования Хафа. Отличие состоит в том, что в пространстве параметров
выбираются два максимума, соответствующие двум прямым. В эксперименте
рассматриваются прямые $y=2x-1$ и $y=-x+101$.

В табл.~\ref{tab:ch3_two_005_nothr} и
\ref{tab:ch3_two_020_nothr} приведены результаты для уровней шума
$\sigma=0.05$ и $\sigma=0.20$ соответственно.

\begin{table}[H]
\centering
\caption{Две прямые: $\sigma=0.05$, без порога}
\label{tab:ch3_two_005_nothr}
\resizebox{\textwidth}{!}{%
\begin{tabular}{|l|l|c|c|c|c|c|c|}
\hline
Конфигурация & Метрика & \texttt{ROW.ROW} & \texttt{COL.ROW} & \texttt{MAX.ROW} & Median & Wiener & \texttt{ROW.ROW}+ESPRIT
\\
\hline
$(-1,101;c=0.5)$ & MSE по $\rho$ & 2448.3 & 1092.5 & 3.16e-04 & 877.4 & 4813.3 & 3.16e-04
\\
\hline
 & MSE по $\theta$ & 0.6169 & 0.6169 & 3.62e-07 & 0.6169 & 0.6169 & 3.62e-07
\\
\hline
 & Активные пиксели & 3216.2 & 4800.8 & 200 & 4993.9 & 9999.8 & 174
\\
\hline
$(2,-1;c=0.8)$ & MSE по $\rho$ & 512.5 & 1906.7 & 1.63e-04 & 1430.9 & 0.3898 & 1.63e-04
\\
\hline
 & MSE по $\theta$ & 0.2150 & 0.2150 & 3.02e-09 & 0.2150 & 0.2150 & 3.02e-09
\\
\hline
 & Активные пиксели & 3216.2 & 4800.8 & 200 & 4993.9 & 9999.8 & 174
\\
\hline
\end{tabular}%
}
\end{table}


\begin{table}[H]
\centering
\caption{Две прямые: $\sigma=0.20$, без порога}
\label{tab:ch3_two_020_nothr}
\resizebox{\textwidth}{!}{%
\begin{tabular}{|l|l|c|c|c|c|c|c|}
\hline
Конфигурация & Метрика & \texttt{ROW.ROW} & \texttt{COL.ROW} & \texttt{MAX.ROW} & Median & Wiener & \texttt{ROW.ROW}+ESPRIT
\\
\hline
$(-1,101;c=0.5)$ & MSE по $\rho$ & 1342.5 & 704.3 & 3.35e-04 & 973.8 & 4817.5 & 3.67e-04
\\
\hline
 & MSE по $\theta$ & 0.6169 & 0.6169 & 3.62e-07 & 0.6169 & 0.6169 & 3.62e-07
\\
\hline
 & Активные пиксели & 4810.1 & 5151.6 & 200 & 5146.7 & 10000 & 133.9
\\
\hline
$(2,-1;c=0.8)$ & MSE по $\rho$ & 1347.8 & 1484.8 & 1.62e-04 & 1394.2 & 0.3266 & 1.58e-04
\\
\hline
 & MSE по $\theta$ & 0.2150 & 0.2150 & 3.02e-09 & 0.2150 & 0.2150 & 3.02e-09
\\
\hline
 & Активные пиксели & 4810.1 & 5151.6 & 200 & 5146.7 & 10000 & 133.9
\\
\hline
\end{tabular}%
}
\end{table}


В случае двух прямых без порога особенно заметен эффект растекания. Метод
\texttt{CSSA COL.ROW} сохраняет несколько тысяч активных пикселей, поэтому
аккумуляторный массив строится не только по точкам двух прямых, но и по
широкому фону вокруг них. В результате максимум, соответствующий истинной
линии, конкурирует с максимумами, созданными дополнительными активными
пикселями. Это объясняет большие ошибки \texttt{COL.ROW} без порога, несмотря
на то, что сам метод сохраняет значимую часть сигнала.

\texttt{MAX.ROW} и \texttt{CSSA ROW.ROW + ESPRIT} ведут себя иначе. Первый
метод жестко ограничивает мощность активного множества числом сохраняемых
максимумов в строках. Второй уменьшает растекание за счет коррекции частот
перед обратным DFT. Поэтому в обеих таблицах эти методы выходят на нижний
уровень, заданный дискретизацией Хафа. Классические фильтры в этой постановке
ведут себя неоднородно: например, Wiener может давать малую ошибку по
$\rho$ для одной из прямых. Однако по $\theta$ ошибка остается большой, так как
после фильтрации активной остается практически вся матрица.

\section{С пороговой фильтрацией}

Теперь рассмотрим результаты после формирования активного множества с помощью
порога. В этих экспериментах основным является порог $1\hat\sigma$; для одной
прямой при $\sigma=0.20$ дополнительно приведен более жесткий порог
$2\hat\sigma$.

\subsection{Одна прямая}

Для одной прямой при пороговой фильтрации оценивается, насколько удаление малых
значений по порогу помогает компенсировать растекание, возникающее после
реконструкции после CSSA.

В табл.~\ref{tab:ch3_one_005_thr1} и \ref{tab:ch3_one_020_thr1} приведены
результаты для порога $1\hat\sigma$ при уровнях шума $\sigma=0.05$ и
$\sigma=0.20$ соответственно.

\begin{table}[H]
\centering
\caption{Одна прямая: $\sigma=0.05$, порог $1\hat\sigma$}
\label{tab:ch3_one_005_thr1}
\resizebox{\textwidth}{!}{%
\begin{tabular}{|l|l|c|c|c|c|c|c|}
\hline
Конфигурация & Метрика & \texttt{ROW.ROW} & \texttt{COL.ROW} & \texttt{MAX.ROW} & Median & Wiener & \texttt{ROW.ROW}+ESPRIT
\\
\hline
$(-1,101;c=0.8)$ & MSE по $\rho$ & 3.16e-04 & 3.16e-04 & 3.16e-04 & 1512.4 & 3.16e-04 & 3.16e-04
\\
\hline
 & MSE по $\theta$ & 1.59e-07 & 1.59e-07 & 1.59e-07 & 0.2248 & 1.59e-07 & 1.59e-07
\\
\hline
 & Активные пиксели & 116 & 100 & 100 & 4.2 & 252.9 & 100
\\
\hline
$(1,0;c=0.8)$ & MSE по $\rho$ & 2.50e-03 & 2.50e-03 & 2.50e-03 & 1437.1 & 2.50e-03 & 2.50e-03
\\
\hline
 & MSE по $\theta$ & 1.43e-06 & 1.43e-06 & 1.43e-06 & 0.4526 & 1.43e-06 & 1.43e-06
\\
\hline
 & Активные пиксели & 116.2 & 100 & 100 & 4.2 & 254.4 & 100
\\
\hline
$(2,-1;c=0.8)$ & MSE по $\rho$ & 0.0233 & 0.0233 & 0.0233 & 4100.9 & 0.0233 & 0.0233
\\
\hline
 & MSE по $\theta$ & 4.22e-06 & 4.22e-06 & 4.22e-06 & 1.5598 & 4.22e-06 & 4.22e-06
\\
\hline
 & Активные пиксели & 99.3 & 51.5 & 100 & 4.3 & 157.2 & 69
\\
\hline
$(2,-40;c=0.8)$ & MSE по $\rho$ & 0.0447 & 0.0447 & 0.0447 & 2553.1 & 0.0447 & 0.0447
\\
\hline
 & MSE по $\theta$ & 4.22e-06 & 4.22e-06 & 4.22e-06 & 1.3702 & 4.22e-06 & 4.22e-06
\\
\hline
 & Активные пиксели & 99.8 & 51.9 & 100 & 3.8 & 159.2 & 69.5
\\
\hline
\end{tabular}%
}
\end{table}


\begin{table}[H]
\centering
\caption{Одна прямая: $\sigma=0.20$, порог $1\hat\sigma$}
\label{tab:ch3_one_020_thr1}
\resizebox{\textwidth}{!}{%
\begin{tabular}{|l|l|c|c|c|c|c|c|}
\hline
Конфигурация & Метрика & \texttt{ROW.ROW} & \texttt{COL.ROW} & \texttt{MAX.ROW} & Median & Wiener & \texttt{ROW.ROW}+ESPRIT
\\
\hline
$(-1,101;c=0.8)$ & MSE по $\rho$ & 3.16e-04 & 3.16e-04 & 3.16e-04 & 3.81e-04 & 3.16e-04 & 3.16e-04
\\
\hline
 & MSE по $\theta$ & 1.59e-07 & 1.59e-07 & 1.59e-07 & 1.59e-07 & 1.59e-07 & 1.59e-07
\\
\hline
 & Активные пиксели & 105.3 & 100.3 & 100 & 103.4 & 287.1 & 78.2
\\
\hline
$(1,0;c=0.8)$ & MSE по $\rho$ & 7.65e-03 & 2.50e-03 & 5.08e-03 & 924.9 & 2.50e-03 & 7.30e-03
\\
\hline
 & MSE по $\theta$ & 1.43e-06 & 1.43e-06 & 1.43e-06 & 0.1296 & 1.43e-06 & 1.43e-06
\\
\hline
 & Активные пиксели & 105.2 & 100 & 100 & 114.7 & 289 & 77.5
\\
\hline
$(2,-1;c=0.8)$ & MSE по $\rho$ & 0.0182 & 0.0232 & 0.0200 & 2628.6 & 394.6 & 0.0194
\\
\hline
 & MSE по $\theta$ & 4.27e-06 & 4.22e-06 & 4.22e-06 & 0.5894 & 0.0408 & 4.31e-06
\\
\hline
 & Активные пиксели & 93.3 & 110.3 & 100 & 92.2 & 266.9 & 57.3
\\
\hline
$(2,-40;c=0.8)$ & MSE по $\rho$ & 0.0279 & 0.0445 & 0.0317 & 2160.3 & 182.6 & 0.0250
\\
\hline
 & MSE по $\theta$ & 4.22e-06 & 4.22e-06 & 4.22e-06 & 1.2487 & 0.0629 & 4.27e-06
\\
\hline
 & Активные пиксели & 93.2 & 110.2 & 100 & 88.9 & 267.3 & 58.2
\\
\hline
\end{tabular}%
}
\end{table}


При пороге $1\hat\sigma$ методы CSSA существенно выигрывают от удаления малых
восстановленных значений. Для \texttt{ROW.ROW} и \texttt{COL.ROW} число
активных пикселей становится сопоставимым с длиной линии, и ошибки переходят
к уровню дискретизации. Это подтверждает, что в варианте без порога основная
проблема заключалась не в потере положения прямой, а в растекании
восстановленного сигнала.

Median ведет себя иначе. Уже при пороге $1\hat\sigma$ для слабого шума он
оставляет только несколько активных пикселей, а при пороге $2\hat\sigma$ и
$\sigma=0.20$ для наклонных прямых активное множество практически исчезает.
Поэтому большие ошибки Median связаны не с ложным фоном, а с потерей самой
линии.

Для проверки чувствительности к более жесткому отсечению дополнительно
рассмотрим порог $2\hat\sigma$ при $\sigma=0.20$.

\begin{table}[H]
\centering
\caption{Одна прямая: $\sigma=0.20$, порог $2\hat\sigma$}
\label{tab:ch3_one_020_thr2}
\resizebox{\textwidth}{!}{%
\begin{tabular}{|l|l|c|c|c|c|c|c|}
\hline
Конфигурация & Метрика & \texttt{ROW.ROW} & \texttt{COL.ROW} & \texttt{MAX.ROW} & Median & Wiener & \texttt{ROW.ROW}+ESPRIT
\\
\hline
$(-1,101;c=0.8)$ & MSE по $\rho$ & 3.16e-04 & 3.16e-04 & 3.16e-04 & 2747.8 & 3.16e-04 & 3.16e-04
\\
\hline
 & MSE по $\theta$ & 1.59e-07 & 1.59e-07 & 1.59e-07 & 0.3934 & 1.59e-07 & 1.59e-07
\\
\hline
 & Активные пиксели & 87.1 & 100 & 100 & 2.8 & 107.8 & 76.3
\\
\hline
$(1,0;c=0.8)$ & MSE по $\rho$ & 8.08e-03 & 2.50e-03 & 5.08e-03 & 1198.7 & 3.33e-03 & 7.38e-03
\\
\hline
 & MSE по $\theta$ & 1.43e-06 & 1.43e-06 & 1.43e-06 & 0.5419 & 1.43e-06 & 1.43e-06
\\
\hline
 & Активные пиксели & 86.6 & 100 & 100 & 2.6 & 106.5 & 75.5
\\
\hline
$(2,-1;c=0.8)$ & MSE по $\rho$ & 0.0193 & 0.0196 & 0.0200 & 18058.4 & 0.0189 & 0.0194
\\
\hline
 & MSE по $\theta$ & 4.31e-06 & 4.76e-06 & 4.22e-06 & 6.4124 & 4.22e-06 & 4.36e-06
\\
\hline
 & Активные пиксели & 65.9 & 34.7 & 100 & 0.1 & 56.8 & 52.9
\\
\hline
$(2,-40;c=0.8)$ & MSE по $\rho$ & 0.0264 & 0.0312 & 0.0317 & 14506.6 & 0.0347 & 0.0250
\\
\hline
 & MSE по $\theta$ & 4.22e-06 & 4.40e-06 & 4.22e-06 & 6.7596 & 4.22e-06 & 4.27e-06
\\
\hline
 & Активные пиксели & 66.5 & 37.7 & 100 & 0.1 & 57.5 & 54
\\
\hline
\end{tabular}%
}
\end{table}


Вариант с порогом $2\hat\sigma$ показывает границу применимости пороговой
фильтрации. Для методов, у которых после восстановления значения на линии
остаются достаточно выраженными, более жесткий порог не ухудшает оценку.
Однако для Median он удаляет почти все активные пиксели, поэтому
преобразование Хафа фактически применяется к неполному изображению линии.
Следовательно, пороговая фильтрация подавляет растекание, но ее нельзя
рассматривать как нейтральный этап: она может как убрать ложные активные
пиксели, так и удалить полезные точки линии.

\subsection{Две прямые}

В табл.~\ref{tab:ch3_two_005_thr1} и \ref{tab:ch3_two_020_thr1} приведены
результаты для порога $1\hat\sigma$ при уровнях шума $\sigma=0.05$ и
$\sigma=0.20$ соответственно.

\begin{table}[H]
\centering
\caption{Две прямые: $\sigma=0.05$, порог $1\hat\sigma$}
\label{tab:ch3_two_005_thr1}
\resizebox{\textwidth}{!}{%
\begin{tabular}{|l|l|c|c|c|c|c|c|}
\hline
Конфигурация & Метрика & \texttt{ROW.ROW} & \texttt{COL.ROW} & \texttt{MAX.ROW} & Median & Wiener & \texttt{ROW.ROW}+ESPRIT
\\
\hline
$(-1,101;c=0.5)$ & MSE по $\rho$ & 3.16e-04 & 3.16e-04 & 3.16e-04 & 1133.5 & 3.16e-04 & 3.16e-04
\\
\hline
 & MSE по $\theta$ & 3.62e-07 & 3.62e-07 & 3.62e-07 & 0.3768 & 3.62e-07 & 3.62e-07
\\
\hline
 & Активные пиксели & 215.9 & 161.9 & 200 & 5.5 & 401.8 & 168.4
\\
\hline
$(2,-1;c=0.8)$ & MSE по $\rho$ & 1.63e-04 & 1.63e-04 & 1.63e-04 & 1741.2 & 1.63e-04 & 1.63e-04
\\
\hline
 & MSE по $\theta$ & 3.02e-09 & 3.02e-09 & 3.02e-09 & 0.6640 & 3.02e-09 & 3.02e-09
\\
\hline
 & Активные пиксели & 215.9 & 161.9 & 200 & 5.5 & 401.8 & 168.4
\\
\hline
\end{tabular}%
}
\end{table}


\begin{table}[H]
\centering
\caption{Две прямые: $\sigma=0.20$, порог $1\hat\sigma$}
\label{tab:ch3_two_020_thr1}
\resizebox{\textwidth}{!}{%
\begin{tabular}{|l|l|c|c|c|c|c|c|}
\hline
Конфигурация & Метрика & \texttt{ROW.ROW} & \texttt{COL.ROW} & \texttt{MAX.ROW} & Median & Wiener & \texttt{ROW.ROW}+ESPRIT
\\
\hline
$(-1,101;c=0.5)$ & MSE по $\rho$ & 3.53e-04 & 3.26e-04 & 3.35e-04 & 1024.4 & 3.33e-04 & 3.62e-04
\\
\hline
 & MSE по $\theta$ & 3.62e-07 & 3.62e-07 & 3.62e-07 & 0.3523 & 3.62e-07 & 3.62e-07
\\
\hline
 & Активные пиксели & 187 & 324.1 & 200 & 119.7 & 359.1 & 114
\\
\hline
$(2,-1;c=0.8)$ & MSE по $\rho$ & 1.60e-04 & 1.63e-04 & 1.62e-04 & 1346.6 & 1.63e-04 & 1.58e-04
\\
\hline
 & MSE по $\theta$ & 3.02e-09 & 3.02e-09 & 3.02e-09 & 0.4939 & 3.02e-09 & 3.02e-09
\\
\hline
 & Активные пиксели & 187 & 324.1 & 200 & 119.7 & 359.1 & 114
\\
\hline
\end{tabular}%
}
\end{table}


После порога $1\hat\sigma$ различия между \texttt{ROW.ROW},
\texttt{COL.ROW}, \texttt{MAX.ROW}, Wiener и
\texttt{CSSA ROW.ROW + ESPRIT} становятся малыми. Для двух прямых это означает,
что порог в основном устраняет лишние активные пиксели, мешавшие выбору двух
максимумов в пространстве Хафа. Median снова остается исключением: число
активных пикселей оказывается слишком малым, поэтому метод теряет часть
геометрической информации о линиях.

\section{Итог сравнения}
Наиболее устойчивыми в рассмотренных конфигурациях оказались методы
\texttt{MAX.ROW} и \texttt{CSSA ROW.ROW + ESPRIT}. В методе
\texttt{MAX.ROW} число активных пикселей контролируется явно: в каждой строке
сохраняется фиксированное число максимумов. В методе
\texttt{CSSA ROW.ROW + ESPRIT} растекание уменьшается за счет коррекции
частот, выполняемой перед обратным DFT. В обоих случаях число активных
пикселей остается небольшим, а ошибки часто совпадают с нижним
уровнем, задаваемым дискретизацией преобразования Хафа. Следовательно,
дальнейшее улучшение результата в таких конфигурациях связано уже не только с
предобработкой, но и с уменьшением шага сетки $(\rho,\theta)$.

Для нескольких прямых сохраняются те же основные закономерности, но возрастает
роль выбора числа компонент и способа выделения нескольких максимумов в
пространстве Хафа. В частности, безопасная оценка ранга для
\texttt{COL.ROW} может быть достаточной для сохранения сигнала, но слишком
большой для эффективного шумоподавления. Поэтому переход к нескольким прямым
требует отдельного анализа ранга, порога активности и процедуры выбора
максимумов в аккумуляторном массиве.
